\section{Service Admission Architecture}
\label{sec:system}
\subsection{Request and Interpretation Contract}
Each request $i$ contains a natural-language description, an arrival time
$a_i$, an absolute deadline $d_i$, an origin, and payload metadata. Times
share the originating controller's clock. Only the
description is sent to the interpreter; the scheduler receives the remaining
metadata through the same interface for every backend. Reference intent
labels and expected service outputs are available exclusively to the
evaluator. Table~\ref{tab:schema} defines the shared interpretation contract.
Its four fields describe the requested service and execution requirements,
with $4\times3\times3\times3=108$ possible tuples. Missing requirements remain
unspecified rather than being inferred from the evaluator's labels.

\begin{table}[t]
\caption{Shared intent contract. OCR denotes optical character recognition.}
\label{tab:schema}
\centering
\begin{tabular}{@{}ll@{}}
\toprule
Field & Allowed values \\
\midrule
Service & Count, detection, OCR, unsupported \\
Locality & Site only, remote allowed, unspecified \\
Quality floor & Standard, high, unspecified \\
Urgency & Normal, urgent, unspecified \\
\bottomrule
\end{tabular}
\end{table}

Jev selects among the allowed values, while the generative backends return
the same fields in a concise structured response. The rule-based interpreter
maps the same input text to this contract. The required meaning of each
output is fixed across interfaces. A request to read an image
locally with high quality and urgent handling, for example, maps to OCR, site
only, high, and urgent. Selecting the execution node remains the scheduler's
responsibility.

\subsection{Validation, Caching, and Admission}
The validator checks that all required fields are present and that each value
belongs to its allowed set. Invalid interpretations do not proceed to
scheduling. Requests also leave the admission path when their deadline or
the available admission capacity is exhausted. This common policy controls
which responses can be executed; semantic correctness must be evaluated
separately against the request's intended meaning.

When caching is enabled, each backend reuses its own validated interpretation
for identical normalized request text under the same extraction policy.
A response becomes reusable only after it has arrived. The cache contains
neither reference labels nor service outputs, and it can retain a
well-formed but semantically incorrect interpretation. Repeated descriptions
may accompany different images, so reusing a decision does not eliminate
the corresponding service execution. All backends use the same cache policy.

Admission and execution share one elapsed-time axis. Waiting for an
interpretation slot and receiving a remote response both consume the budget
between $a_i$ and $d_i$. A call that outlives its request's deadline continues
to occupy its slot until the call terminates, while new requests continue
to arrive. Interpretation latency thus affects both service feasibility
and queue occupancy.

\subsection{Shared Scheduling and Execution}
After interpretation, the scheduler reads the current worker state. Only an
explicit remote permission allows a payload to leave its originating site;
unspecified locality defaults to local execution. A high quality requirement
selects the high tier, and other requests use the standard tier. Urgent
requests take priority over ordinary requests without interrupting active
work. These shared defaults can make distinct field values operationally
equivalent, as with normal and unspecified urgency.

For the real service, let $t_i$ denote the time at which request $i$ is
scheduled, after its interpretation has completed. Its selected tier is
$q_i$, and its priority is $p_i$, with smaller values indicating higher
priority. Let $\mathcal{N}_i$ contain the nodes that provide the interpreted
service, satisfy its placement restriction, and have queue capacity.
For a node $n\in\mathcal{N}_i$, let $b_n$ be the estimated finish time of
active work, or zero if idle, and let $h_n(q)$ be the calibrated duration of
serving tier $q$, including transfer and response overhead. Finally,
$\mathcal{Q}_n^{\leq p_i}$ is the set of pending jobs with equal or higher
priority than request $i$, and $q_j$ is pending job $j$'s tier. All node
quantities refer to the state observed at $t_i$. The predicted finish time
$\widehat f_{in}$ is
\begin{equation}
\widehat f_{in}=\max(t_i,b_n)
  +\sum_{j\in\mathcal{Q}_n^{\leq p_i}}h_n(q_j)+h_n(q_i).
\label{eq:finish}
\end{equation}
The scheduler selects the node with the smallest predicted finish time among
those satisfying $\widehat f_{in}\leq d_i$, breaking ties in favor of the
local node. It rejects the request if no such node exists. Equation~\ref{eq:finish}
is a shared admission estimate: execution-time variation and subsequent
higher-priority arrivals can invalidate its prediction. Correct completion
therefore depends on the observed result and finish time, not merely on
passing the admission test.

An accepted request transfers its payload to the selected worker and returns
the service output to the originating controller. Payloads are transferred
only after node selection. The evaluator checks the actual service, node,
tier, priority, output, and deadline against the reference requirements.
All elapsed times are measured at the originating controller. Locality
constraints govern service payloads; the textual description still reaches
the chosen interpretation API. The scheduler and these checks are common
to all backends, allowing the comparison to focus on interpreter substitution.
